\documentclass[11pt]{article}
\usepackage{amsmath,amssymb,amsthm}
\newtheorem{lemma}{Lemma}
\newtheorem{proposition}{Proposition}
\theoremstyle{remark}\newtheorem*{remark}{Remark}
\title{Deferred readouts in the two-stage bit interchange}
\date{Draft, October 9, 2026}
\begin{document}
\sloppy
\maketitle

This note describes the round-seven bit side. It keeps the round-six interchange (flag basis, retained point totals
with copied centres, side roles batched one level down, data-entrance run) and changes three things: the frames of
the side program are lifted and its fan-out copies are made late, and the stage-one word is reordered so that most
slots read their garbage out at a nonzero frame. We give two witnesses with two different side programs: witness~1
uses PR~\#62's producer (\S\ref{sec:pr62}), witness~2, the headline, PR~\#57's joint frame compiler
(\S\ref{sec:jf}); the second needs a stricter rule for the first phase of the reordered word (\S\ref{sec:closure}).
All objects are frozen in \texttt{certificates/round7/}; \texttt{scripts/certificate\_round7.py} rebuilds both
child-width histograms from the schedules, and the checks in \texttt{independent/deferred-readout/} (witness~1) and
\texttt{independent/joint-frame-stack/} (witness~2) verify every fact used below. Throughout, $h=23$,
$m=h^2$, $v=\binom h3$, $N=v^2$, $G=I-J/9$, and for a triple $T$ the root frame of the output wire $y_T$ is
$t_T^\perp=\ker\,(9t_T-3\mathbf 1)^{\mathsf T}$; the frame of a retained total of $c$ is $U_c=\ker(\mathbf 1-3e_c)^{\mathsf T}$.

\section{Witness 1: PR \#62's producer}\label{sec:pr62}
The side DAG has the $v$ triples as leaves and disjoint additions as nodes; its roots are the $3v$ outputs
$\mathrm{out}(c,T)=\{S: S\cap T=\{c\}\}$ and the $h$ retained totals $\mathrm{star}(c)$. We build it with PR~\#62's
producer (ikeboy): interval strips and core-aware pair assembly, with PR~\#41's alternating order and links
(RaD/hipotures), and links chosen by a moment-weighted matching in the spirit of PR~\#44 (rohanarun). A link hands the
non-pivot slot of a donor gate to a later use of the same value, so the number of slots is
$R=\#\text{additions}+3v+h-\#\text{links}=28866$. Only the DAG enters the certificate, and it is checked from the
definitions: every addition is disjoint, every output and total has the support above, and every node is reachable
from a root (\texttt{check\_lifted.py}, part~A; \texttt{check\_word.py}, part~A).

\section{Witness 2: the joint frame compiler}\label{sec:jf}
Witness~2 builds the side DAG with the same producer and PR~\#69's balanced coarse sums (eumemic; every four-edge
coarse sum as two column pair sums and a total) and compiles it with PR~\#57's joint frame compiler (eumemic) in
PR~\#60's rank-first reclamation order (chafreaky): roles are reused across regions of the DAG, and every gate acts at
the frame of its region. The compiled program has $R=27794$ roles. Its regions are lifted as below
($M_g=\operatorname{span}(g)+(U_g\cap F_{j(g)})$ for a region $g$, with $U_g$ the common kernel of the root covectors
reachable from $g$ along the roles), and late copies sit at exact intersections of their use frames. The compiler's
region maps are $\mathbb F_2$-linear, so its word identity holds over $\mathbb F_2$, which is all the bit network
needs (\S\ref{sec:uses}); it is checked by an exact symbolic $\mathbb F_2$ replay with every scratch bit a separate
symbol, for the compiled, the lifted and the final program.

\section{Lifted frames and late copies}
For an addition $n$ let $U_n$ be the common kernel of the root covectors reachable from $n$ and $F_j$ the span of
the first $j$ rows of a fixed integer flag matrix. The gate of $n$ acts at
\[ M_n=\operatorname{span}(n)+(U_n\cap F_{j(n)}), \]
with a lift index $j(n)$ per addition. This is Paureel's complement-frame construction: it is valid for any frames
with (H1) leaf frames $\langle t_S\rangle$, (H2) frames nested along every slot (DAG edges, links, copies),
(H3) roots inside $t_T^\perp$ or $U_c$, (H4) $G$ nondegenerate on every frame. Nesting is a theorem given
$j(n)\le j(t)$ on every slot successor pair, $\operatorname{span}(x)\le\operatorname{span}(t)$ on every link and
$\operatorname{span}(n)$ inside every root it reaches directly; then $U_n\le U_t$ and $M_n\le M_t$. In this witness
every lifted frame equals $U_n$. Dimensions are exact (modular law, ranks with Hadamard certificates or Bareiss).

\paragraph{Late copies.} A node with several free uses is copied for each of them. A copy made early sits at $M_n$;
a \emph{late} copy for the user $\ell_i$ of an ordered list $\ell_1,\dots,\ell_k$ is made when the pivot has climbed to
\[ C_i=\textstyle\bigcap_{i'\ge i} M(\ell_{i'}), \qquad M_n\le C_1\le\dots\le C_{k-1}\le M(\ell_k), \]
and starts there, so the copy saves the rank $\dim C_i-\dim M_n$. Each copy is placed as late as possible (before its
user's gate, the pivot's next gate and the next copy). The intersections are needed: without them consecutive pivot
frames are not nested (1396 of 1396), and a copy moved after the pivot's next gate breaks the replay.

\section{The B-defer word}
In stage one each slot $u$ starts with its garbage $z_u$, and the word reads it out of every target it reaches:
$y_T\mathrel{-}=C_{T,u}z_u$, where $C=JL$ is given by the transpose sweep of $L$. In rounds five and six all these
readouts happen at frame $0$, before the second pass. B-defer splits the second pass in two phases.

\begin{itemize}
\item \emph{Phase 1} is the closure of the retained totals' last operations under per-slot sequence precedence: every
  op that must precede a total's completion (for witness~2 with the frame-order edges of \S\ref{sec:closure}). The totals complete, and their copies read every target containing
  their centre at $0$ (copied centres, as in round six).
\item A slot that phase 1 neither reads nor writes still holds its garbage after phase 1. It reads it out
  \emph{now}, at a frame $\sigma_u\le F_0(u)\cap\bigcap_{T}t_T^\perp$ (over its targets $T$), then the rest of $L$ runs,
  then the output reads at $t_T^\perp$, then $L^{-1}$ and $V^{-1}$ exactly mirror it.
\end{itemize}
Only the topological order of the word changes, so it is still a mirror and still restores arbitrary scratch. The
slot's chain now starts at $\sigma_u$: its corner has rank $r_u=h-\dim\sigma_u$ and its auxiliary edge rank $m-r_u$.

\paragraph{The sorting rule.} On a target $y_T$ the readout frames, in time order, are $0$ (non-deferred slots and
centre copies), the deferred $\sigma_u$, and $t_T^\perp$; they must form a chain. The word reads the deferred slots
in increasing $\dim\sigma_u$, and applies the deferred $V$ gates (below) in increasing $\dim F_0$. With this rule
every $y_T$ sequence is nested in time order, exactly over $\mathbb Q$, for both the $\mathbb Z$ and the
$\mathbb F_2$ supports of $C$; the reversed order breaks 11879 chains.

\section{The phase-1 closure}\label{sec:closure}
Phase~1 must contain every op that has to run before a retained total completes. Read/write precedence alone is
enough for witness~1, whose slots move along a single chain of node frames. Under the joint frame compiler a role
passes through several regions, and an op that touches a role at an earlier frame of its chain must run before any
op that touches it at a later frame, or the role would have to step down its chain. The rule is therefore the
closure of the retained totals' last writes under
\begin{itemize}
\item read/write edges: an op follows the last write of each role it reads or writes, and a write follows the
  earlier reads of that role;
\item frame-order edges: an op touching role $b$ at chain position $k$ follows every op touching $b$ at a lower
  position (positions from the walk of the chain in the original op order, frames compared as exact subspaces).
\end{itemize}
For witness~2 the read/write closure has 25942 ops and the full closure 31542. The checks walk the actual reordered
op sequence (phase~1, centre reads, deferred readouts, deferred $V$ gates, phase~2, output reads), and require every
event to sit at a frame of both roles' chains with pointers that only move up; with the read/write closure alone that
walk fails. The reordering keeps the version of every role seen by every op, so $L$ is unchanged as a map.

\section{V leaves as late copies on $X_S$}
A leaf $S$ is used by several gates. Instead of copying $x_S$ by fan-out, every use $i$ gets its own $V$ gate on the
data wire $X_S$, into a slot that starts at
\[ s_i=N_i\cap N_{i+1}\cap\cdots \qquad(\text{falling back to }\langle t_S\rangle\text{ if $G$ is degenerate on it}), \]
where $N_j$ is the first frame of use $j$. The $V$ gates are then late copies on $X_S$: its frame chain
$\langle t_S\rangle\le s_{i_1}\le s_{i_2}\le\dots\le F$ has total rank $h-1$, as before, and the use slots start at
$s_i$ instead of $\langle t_S\rangle$. The $V$ gates of deferred slots come after their readouts; in time the $X_S$ chain
is $\langle t_S\rangle$, the early $V$ frames (by dimension), the deferred $V$ frames (by the rule above), $F$, and it
is checked nested exactly in that order. The number of slots is unchanged.

\section{What the certificate uses}\label{sec:uses}
\paragraph{Only the $\mathbb F_2$ flag identity.} The bit network is used over $\mathbb F_2$: the word must add exactly
$x_T$ to every $y_T$ and restore every slot, for arbitrary scratch. The replay checks this over $\mathbb F_2$ with
the parity coefficients (the actual circuit) and, separately, over $\mathbb Z$ with the exact path-count
coefficients (witness~1); the histogram and $a_b$ are the same with either support. For witness~2 the replay is
symbolic over $\mathbb F_2$.

\paragraph{The stage-two complement.} Stage two is the complement time-reversal of stage one, with the same edge
multiset. This is the same assumption as in round six; it is stated, not machine-checked here.

\paragraph{Frames and genericity.} Every $\sigma_u$ and $s_i$ is a primitive integer basis (entries at most 14).
Exactly over $\mathbb Q$: $\sigma_u\le F_0(u)$, $\sigma_u\le t_T^\perp$, $t_S\le s_i\le N_i$, $G$ nondegenerate. Every
high-rank step ($2r>h$) that the reordered word adds (slot starts, gauged corners $\sigma_u\to F$, $y_T$ and $X_S$
chain steps) has the generic north-east pivots under both conjugations at the integer point of
\texttt{check\_lifted.py}, so the side lemma batches it.

\section{Accounting}
Per stage and invocation: every slot contributes its auxiliary edge (block $m-2r_u$) and its corner, both batched
by the inner lemma, and its chain steps; every centre copy its step $U_c\to0$ of rank $h-1$; every $y_T$ and every
$X_S$ its chain of total rank $h-1$; the data entrance its block $m-4h+2$ and corner; one copy-correction singleton per
pair. The rank sum is exactly $s=Wm-N+L$ with $W=2N+2vR$, $L=2vh(h-1)$, so the histogram re-partitions the round-six
budget. The moment certificate gives, with the staircase entrance corner (runs $h-2,h-5,1^6$) and with the round-six
corner,
\begin{align*}
 \text{witness 2:}\quad a_b&=32761/(5\cdot10^8), &a_b&=12899/(2\cdot10^8),\\
 \text{witness 1:}\quad a_b&=31987/(5\cdot10^8), &a_b&=12599/(2\cdot10^8),
\end{align*}
and with round six's complex interchange ($a_c=36926111/(5\cdot10^{11})$, $h=24$) the assembly gives the headline
$\kappa=3275885357429/(5\cdot10^{16})\approx6.55177\times10^{-5}$ (witness~2, staircase corner); witness~2 with the
round-six corner gives $3224542033151/(5\cdot10^{16})$, and witness~1 gives $1599247689723/(2.5\cdot10^{16})$ and
$6299103187973/10^{17}$. The staircase corner is certified by \texttt{check\_stair.py} at the same integer point;
the entrance edge is not changed by the reordering.

\end{document}
